\documentclass[12pt]{article}

% === Packages ===
\usepackage[margin=1in]{geometry}
\usepackage{amsmath, amssymb}
\usepackage{setspace}
\usepackage{titlesec}
\usepackage{enumitem}
\usepackage{helvet}
\renewcommand{\familydefault}{\sfdefault} % modern sans-serif look

% === Formatting tweaks ===
\setstretch{1.15}
\setlist[enumerate]{itemsep=0.4em, topsep=0.2em}
\titleformat{\section}{\Large\bfseries\sffamily}{\thesection.}{1em}{}
\titleformat{\subsection}{\large\bfseries\sffamily}{\thesubsection}{1em}{}


% === Title ===
\title{\vspace{-2em}\textbf{Applied Mathematics Oral Exam Questions and Model Answers}\\
\large IBFE}
\author{Prepared for Graduate Oral Examination}
\date{}
\begin{document}
\maketitle
% =========================
% IBFE oral exam questions (from IBFE Core Ideas tutorial)
% =========================

\begin{enumerate}

\item \textbf{Kinematics: What is the configuration map, and what are $F$ and $J$?}\\
\textbf{Sample answer:} The structure is described in reference coordinates $X\in B_0$, and its current position is
$x=\chi(X,t)$. The deformation gradient is $F_\chi(X,t)=\nabla_X \chi(X,t)$, which maps reference line elements via
$d x = F_\chi\, dX$. The Jacobian $J(X,t)=\det(F_\chi)$ maps volumes: $dv = J\, dV$.

\item \textbf{Stress: What is the first Piola--Kirchhoff stress, and how is it defined from hyperelasticity?}\\
\textbf{Sample answer:} For a hyperelastic material we choose a strain energy density $W(F_\chi)$ per reference volume.
The first Piola--Kirchhoff stress is $P(X,t) = \dfrac{\partial W}{\partial F_\chi}(F_\chi(X,t))$. It is the stress measure
that is naturally work-conjugate to $F_\chi$ in a reference description.

\item \textbf{Traction: What does $P$ \emph{mean} physically?}\\
\textbf{Sample answer:} The most direct physical meaning is traction on a \emph{reference} surface patch:
if $N(X)$ is the outward unit normal on $\partial B_0$, then the reference traction is
$t_0(X,t)=P(X,t)\,N(X)$. It is force per reference area, with direction in spatial coordinates.

\item \textbf{Two-point tensor: Why do people call $P$ a ``two-point'' tensor? Give the component form.}\\
\textbf{Sample answer:} $P$ maps a \emph{reference} normal to a \emph{spatial} traction, so it mixes material and spatial
information. In components, $(t_0)_i = P_{iI}\,N_I$, where $i$ is a spatial index and $I$ is a material index. So $P$ is
a mixed (spatial $\times$ material) tensor.

\item \textbf{Reference momentum balance: State the strong form on $B_0$ and interpret each term.}\\
\textbf{Sample answer:} The reference strong form is
$\rho_0(X)\,\ddot{\chi}(X,t)=\nabla_X\cdot P(X,t) + B(X,t)$ in $B_0$.
Here $\rho_0$ is reference density, $\nabla_X\cdot P$ is internal force per reference volume, and $B$ is a body force per
reference volume (e.g.\ gravity or other prescribed forcing).

\item \textbf{Divergence theorem step: How do you get from boundary traction to $\nabla_X\cdot P$?}\\
\textbf{Sample answer:} For any material subregion $V_0\subset B_0$,
the net contact force is $\int_{\partial V_0} P N\, dA$. By the reference divergence theorem,
$\int_{\partial V_0} PN\, dA = \int_{V_0} \nabla_X\cdot P \, dV$. This shows $\nabla_X\cdot P$ is force per unit reference volume.

\item \textbf{Weak form derivation: Starting from strong form, derive the weak form (structure side).}\\
\textbf{Sample answer:} Multiply the strong form by a smooth test function $\psi(X)$ and integrate over $B_0$:
$\int_{B_0}\rho_0 \ddot{\chi}\cdot \psi\, dV = \int_{B_0} (\nabla_X\cdot P)\cdot \psi\, dV + \int_{B_0} B\cdot \psi\, dV$.
Then use the identity $\nabla_X\cdot(P^T\psi)=(\nabla_X\cdot P)\cdot\psi + P:\nabla_X\psi$ and integrate by parts to obtain
\[
\int_{B_0} (\nabla_X\cdot P)\cdot \psi\, dV
= -\int_{B_0} P:\nabla_X\psi\, dV + \int_{\partial B_0} (PN)\cdot \psi\, dA.
\]
Substitute back to get the weak form, including any boundary traction term.

\item \textbf{Boundary term: When does the boundary traction term vanish, and what does that assumption mean physically?}\\
\textbf{Sample answer:} It vanishes if the boundary is traction-free ($PN=0$ on $\partial B_0$) or if the material has no
boundary (closed manifold, so $\partial B_0=\emptyset$). Physically, that means no external contact traction is applied on the
solid boundary in the reference description (or there is no boundary).

\item \textbf{Defining the Lagrangian force density $F$: What does ``define by duality'' mean here?}\\
\textbf{Sample answer:} Instead of computing $\nabla_X\cdot P$ pointwise, we define a force density $F(X,t)$ so that for all
test functions $\psi$,
\[
\int_{B_0} F\cdot \psi\, dV
= -\int_{B_0} P:\nabla_X\psi\, dV + \int_{B_0} B\cdot \psi\, dV
\quad (+\text{boundary work if present}).
\]
This is exactly the virtual-work statement: $F$ is the force density whose action on $\psi$ equals the internal virtual work.

\item \textbf{Coupling: Write the two IB coupling equations (spreading and velocity interpolation) and interpret them.}\\
\textbf{Sample answer:} Spreading: $f(x,t)=\int_{B_0} F(X,t)\,\delta_h(x-\chi(X,t))\, dV$ gives the Eulerian force density
in Navier--Stokes. Interpolation: $\partial_t \chi(X,t)=\int_\Omega u(x,t)\,\delta_h(x-\chi(X,t))\, dx$ moves the structure
with a smoothed no-slip/no-penetration condition. $\delta_h$ is a regularized delta kernel tied to grid spacing.

\item \textbf{Active stress: How do you add active tension in IBFE?}\\
\textbf{Sample answer:} Use an additive split $P=P_p+P_a$. $P_p$ comes from $W(F_\chi)$ and $P_a$ is prescribed by an
active model. Both enter identically in the weak form through $-\int P:\nabla_X\psi\, dV$, so implementation is: evaluate
$P_p$ and $P_a$ at quadrature points, add them, and assemble forces as usual.

\item \textbf{Fiber active model: Explain the form $P_a = J\,T\,a\otimes f_0$ and define all symbols.}\\
\textbf{Sample answer:} Choose a unit reference fiber direction field $f_0(X)$ and a scalar tension $T(X,t)\ge 0$.
Let $F_\chi=\nabla_X\chi$, $J=\det(F_\chi)$, and define the deformed fiber direction (not normalized) $a(X,t)=F_\chi f_0$.
Then a simple PK1 active stress is $P_a = J\,T\,a\otimes f_0$. This represents tension generated along fibers and is written
in reference variables, so it drops directly into the weak form.

\item \textbf{Worked example: Uniaxial stretch. Given $\chi(X)=(\lambda X_1, X_2, X_3)$, compute $F_\chi$, $J$, $C$, $E$.}\\
\textbf{Sample answer:} $F_\chi=\mathrm{diag}(\lambda,1,1)$, $J=\lambda$.
$C=F_\chi^T F_\chi=\mathrm{diag}(\lambda^2,1,1)$.
$E=\tfrac12(C-I)=\tfrac12\,\mathrm{diag}(\lambda^2-1,0,0)$.

\item \textbf{Constitutive example: For compressible neo-Hookean $W=\frac{\mu}{2}(I_1-3)-\mu\ln J +\frac{\kappa}{2}(\ln J)^2$,
what is a standard closed form for $P$?}\\
\textbf{Sample answer:} A common formula is
$P=\mu\left(F_\chi - F_\chi^{-T}\right) + \kappa(\ln J)\,F_\chi^{-T}$.
For $F_\chi=\mathrm{diag}(\lambda,1,1)$, $F_\chi^{-T}=\mathrm{diag}(1/\lambda,1,1)$ and $\ln J=\ln\lambda$, so you can write
$P$ componentwise and see that lateral components generally are nonzero under a prescribed kinematic stretch.

\item \textbf{Concept check: Why does a prescribed uniaxial \emph{stretch} generally produce nonzero lateral stresses?}\\
\textbf{Sample answer:} Because holding $\chi_2=X_2$ and $\chi_3=X_3$ constrains lateral deformation to be zero, which is
not the same as a uniaxial-stress test. A true uniaxial-stress test would enforce $\sigma_{22}=\sigma_{33}=0$ and solve for
lateral contraction (via Poisson-type response).



\end{enumerate}
\end{document}